\documentclass[10pt,a4paper]{amsart}
\usepackage[margin=19mm]{geometry}
\usepackage{amsmath,amssymb,mathtools}
\usepackage[scr=rsfs,cal=euler]{mathalfa}
\usepackage{xfrac}
\usepackage[unicode,hidelinks,bookmarks=false]{hyperref}
\newcommand{\ie}{i.e.\ }
\newcommand{\cf}{cf.\ }
\newcommand{\resp}{respectively\ }
\newcommand{\Subsection}{\subsection}
\providecommand{\nobreakdash}{\nobreak\hbox{-}}
\setlength{\emergencystretch}{2em}
\multlinegap0pt
\usepackage{amssymb}
\usepackage{bm}
\usepackage{mathtools}
\usepackage[shortlabels]{enumitem}
\newtheorem{definition}{Definition}
\newtheorem{theorem}{Theorem}
\title{The application of decay character on the global behavior of damped wave equation with Riesz potential-type power nonlinearity}
\author{Trung Loc Tang, Dinh Van Duong, Duc An Phan}
\date{Source: arXiv:2602.16293v3 (CC BY 4.0)}
\begin{document}
\maketitle

\noindent\emph{Source and licence.} The application of decay character on the global behavior of damped wave equation with Riesz potential-type power nonlinearity. Version arXiv:2602.16293v3 (CC BY 4.0), distributed by arXiv under the Creative Commons Attribution 4.0 International licence, http://creativecommons.org/licenses/by/4.0/, as recorded in the arXiv licence metadata for this exact version and shown on the abstract page https://arxiv.org/abs/2602.16293v3. Full licence notice: \texttt{licenses/arxiv-2602.16293-CC-BY-4.0.txt}.\par\smallskip

\noindent\textit{Editorial scope.} Counted unit: the article's theorem Thm\_Blow\_up (source0.tex lines 407--417). The article's complete original proof of it is delivered with it (source0.tex lines 434--527, one \texttt{proof} environment of 94 lines, which the article places away from the statement). No mathematics is written, repaired or re-derived: the statement and the proof are the article's own lines, in the article's own notation, and no retained line is substituted or re-flowed.  Retained uncounted as the source-local context the proof needs: lines 74--79; lines 399--405; lines 401--403.  Nothing was omitted from any retained span, so the package carries no omission record.  No retained line contains a citation, so the wrapper carries no bibliography and no bibliography entry.  No retained line contains a figure environment, an image-inclusion command or drawing code, so no image is delivered and the package declares no asset dependency.  Editorial formatting only, in addition: an AMS \texttt{amsart} preamble that restates the article's own declarations and macros used by the retained lines, namely the environments \texttt{definition}, \texttt{theorem} on separate counters, together with the standard packages those lines need.  The retained environments print in this document as Definition 1, Theorem 1 in order of appearance, whereas the article prints the same items with the numbering of its own full document.  The complete native source of the article as distributed in the arXiv e-print for this exact version is bundled unchanged as \texttt{sources/194867/source0.tex}.
\begin{align}\label{Semilinear_Damped_Waves}
\begin{cases}
u_{tt}(t,x)-\Delta u(t,x)+u_t(t,x)=\mathcal{I}_\gamma(|u|^p)(t,x), &x\in \mathbb{R}^n, t>0,\\
u(0,x)=\varepsilon u_0(x),\ u_t(0,x)=\varepsilon u_1(x),&x\in \mathbb{R}^n,
\end{cases}
\end{align}

\begin{definition}[\textbf{Weak solution}] \label{weak_def}
Let $p > 1$ and $T \in (0, \infty)$. We say that $u \in \mathcal{C}([0,T), L^2)$ is a local (in time) weak solution to \eqref{Semilinear_Damped_Waves} if  for any test functions $\chi \in \mathcal{C}_0^\infty([0,T))$ and $\phi \in \mathcal{C}_0^{\infty}(\mathbb{R}^n)$,  the following relation holds: 
\begin{align}
\int_0^{T}\int_{ \mathbb{R}^n}\mathcal{I}_{\gamma} (|u|^p)(t,x)\phi(x)\chi(t)dxdt = \int_0^{T}\int_{ \mathbb{R}^n}(\partial_t^2-\Delta+\partial_t)u(t,x)\phi(x)\chi(t)dxdt. \label{eqweak}
\end{align}
If $T = \infty$, then we say that $u$ is a global (in time) weak solution to \eqref{Semilinear_Damped_Waves}.
\end{definition}

\begin{align}
\int_0^{T}\int_{ \mathbb{R}^n}\mathcal{I}_{\gamma} (|u|^p)(t,x)\phi(x)\chi(t)dxdt = \int_0^{T}\int_{ \mathbb{R}^n}(\partial_t^2-\Delta+\partial_t)u(t,x)\phi(x)\chi(t)dxdt. 
\end{align}

\begin{theorem}[\textbf{Blow-up}] \label{Thm_Blow_up}
Let $n \geq 1$, $q \in (0, n/2)$ and $\gamma\in [0,n)$. The exponent $p$ fulfills
\begin{align}
1<p<p_{\mathrm{crit}}(n, q, \gamma). \label{condition-Blowup}
\end{align}
Let us assume $\left(u_0,u_1\right)\in \mathcal{Y}^q\times\mathcal{Y}^q$ such that
\begin{equation} \label{assumption_initial_data}
u_0(x)+u_1(x)\gtrsim\langle x\rangle^{-n+q},
\end{equation}
where we denote $\langle x\rangle:= \left(1+|x|^2\right)^{1/2}$ by the Japanese bracket for $x \in \mathbb{R}^n$. Then, there is no global (in time) weak solution by mean of Definition \ref{weak_def} to \eqref{Semilinear_Damped_Waves}.
\end{theorem}

\begin{proof}[\textbf{Proof of Theorem \ref{Thm_Blow_up}.}]
At first, let us introduce the test functions $\chi= \chi(t)$ and $\phi=\phi(x)$ satisfying the following properties:
\begin{align*}
&1.\quad \chi \in \mathcal{C}_0^\infty([0,\infty)) \text{ and }
\chi(t)=
\begin{cases}
1 &\text{ if } 0 \leq t \leq 1/2, \\
\text{decreasing} &\text{ if } 1/2\leq t\leq 1, \\
0 &\text{ if } t \geq 1,
\end{cases} & \nonumber \\
&2.\quad \phi \in \mathcal{C}_0^\infty(\mathbb{R}^n) \text{ and }
\phi(x)=
\begin{cases}
1 &\text{ if }0 \leq |x| \leq 1/2, \\
\text{decreasing} &\text{ if } 1/2\leq |x| \leq 1, \\
0 &\text{ if } |x| \geq 1,
\end{cases} & \nonumber \\
&3.\quad \chi^{-\frac{p'}{p}}(t)\big(|\chi'(t)|^{p'}+|\chi''(t)|^{p'}\big) \text{ and } \phi^{-\frac{p'}{p}}(x)\left|\Delta \phi(x)\right|^{p'} \text{ are bounded }.
\end{align*}
where $p'$ stands for the conjugate of $p$. For $R>0$, we define $\phi_R(x):=\phi(R^{-1}x)$ and $\chi_R(t):=\chi(R^{-2}t)$. Assume by contradiction that $u=u(t,x)$ is a global (in time) weak solution to \eqref{Semilinear_Damped_Waves}.
 We apply Definition \ref{weak_def} with $\chi$ and $\phi$ replaced by $\chi_R$ and $\phi_R$, respectively.
From this, by integration by parts, the relation \eqref{eqweak} implies that
\begin{align}\label{eqintegral}
&-\varepsilon\int_{\mathbb{R}^n}(u_0(x)+u_1(x))\phi_R(x)dx+\int_0^{\infty}\int_{\mathbb{R}^n}u(t, x)\phi_R(x)\partial_t^2\chi_R(t) dxdt \notag\\
&-\int_0^{\infty}\int_{\mathbb{R}^n} u(t, x)\phi_R(x)\partial_t\chi_R(t)dxdt-\int_0^{\infty}\int_{\mathbb{R}^n}u(t, x)\Delta  \phi_R(x)\chi_R(t)dxdt \notag\\ =&\int_0^{\infty}\int_{\mathbb{R}^n}\mathcal{I}_{\gamma} \left(|u|^p\right)(t,x)\phi_R(x)\chi_{R}(t)dxdt.
\end{align}
We define
$$
\mathcal{J}(R):=\int_0^{\infty}\int_{\mathbb{R}^n}|u(t, x)|^p \phi_R(x)\chi_R(t)dxdt=\int_0^{R^2}\int_{0 \leq|x| \leq R}|u(t, x)|^p \phi_R(x)\chi_R(t)dxdt.
$$
We deduce from \eqref{eqintegral} the inequality
\begin{align}
&\int_0^{R^2}\int_{0 \leq |x| \leq R} \mathcal{I}_{\gamma} \left(|u|^p\right)(t,x)\phi_R(x)\chi_R(t)dxdt+\varepsilon\int_{0 \leq |x| \leq R}(u_0(x)+u_1(x))\phi_R(x)dx \notag\\
\leq & \int_{R^2/2}^{R^2}\int_{0 \leq |x| \leq R} |u(t, x)| |\phi_R(x) \partial_t\chi_R(t)|dxdt+\int_{R^2/2}^{R^2}\int_{0 \leq |x| \leq R} |u(t, x)||\phi_R(x) \partial_t^2\chi_R(t)|dxdt \notag\\
&+
\int_0^{R^2}\int_{R/2 \leq |x| \leq R}|u(t, x)||\Delta \phi_R(x)|\chi_R(t)dxdt \notag\\
=:&\,\,\mathcal{J}_{1,R} + \mathcal{J}_{2,R} + \mathcal{J}_{3,R}. \label{eqabs}
\end{align}
Using H\"older's inequality with $1/p + 1/p' = 1$ and the change of
variables $\overline{x}=R^{-1}x$, $\overline{t}=R^{-2}t$, we obtain
\begin{align}
\mathcal{J}_{1,R}
\lesssim & \mathcal{J}_R^{\frac{1}{p
}}\left(\int_{R^2/2}^{R^2}\int_{0 \leq |x| \leq R}\chi^{-\frac{p'}{p}}_R(t) |\partial_t\chi_R(t)|^{p'}\phi_R(x)dxdt\right)^{\frac{1}{p'}} \notag\\
\lesssim& \mathcal{J}_R^{\frac{1}{p}}R^{-2+\frac{n+2}{p'}} \left(\int_{1/2}^{1}\int_{0 \leq |\bar{x}| \leq 1} |\chi^{\prime}(\overline{t})|^{p'}|\chi(\overline{t})|^{-\frac{p^{\prime}}{p}}\phi(\overline{x}) d\overline{x}d\overline{t}\right)^{\frac{1}{p'}} \sim \mathcal{J}_R^{\frac{1}{p}}R^{-2+\frac{n+2}{p'}}.\label{eqHolder1}
\end{align}
The same argument yields
\begin{equation}\label{eqHolder2}
\mathcal{J}_{2,R} \lesssim \mathcal{J}_R^{\frac{1}{p}}R^{-4+\frac{n+2}{p'}}.
\end{equation}
For the term $\mathcal{J}_{3,R}$, we treat it as follows:
\begin{align}
\mathcal{J}_{3,R}
&\lesssim \mathcal{J}_R^{\frac{1}{p}}
\left(\int_0^{R^2}\int_{R/2 \leq |x| \leq R} |\Delta\phi_R(x)|^{p'} \phi_R^{-\frac{p'}{p}}(x)\chi_R(t)dxdt\right)^{\frac{1}{p'}} \lesssim \mathcal{J}_R^{\frac{1}{p}}R^{-2+\frac{n+2}{p'}}. \label{eqHolder3}
\end{align}
By \eqref{eqabs}, \eqref{eqHolder1}, \eqref{eqHolder2} and \eqref{eqHolder3}, it follows that
\begin{align} \label{I_u^p}
\int_0^{R^2}\int_{0 \leq |x| \leq R} \mathcal{I}_{\gamma} \left(|u|^p\right)(t,x)\phi_R(x)\chi_R(t)dxdt &+\varepsilon\int_{0 \leq |x| \leq R}(u_0(x)+u_1(x))\phi_{R}(x)dx\notag\\ &\quad\leq C \mathcal{J}_R^{\frac{1}{p}}R^{-2+\frac{n+2}{p'}}.
\end{align}
To control the first quantity on the left-hand side of \eqref{I_u^p}, we begin with the representation of the Riesz potential
\begin{align*}
\mathcal{I}_{\gamma}\left(|u|^p\right)(t,x)= C_\gamma\int_{\mathbb{R}^n} \displaystyle\frac{|u(t, y)|^p}{|x-y|^{n-\gamma}} d y
\geq & \,\,C_\gamma\int_{\mathcal{D}_R} \displaystyle\frac{|u(t, y)|^p}{|x-y|^{n-\gamma}} d y
\end{align*}
where $ \mathcal{D}_R=\left\{y \in \mathbb{R}^n, |y| \leq R\right\}$.
Now observe that for $x \in \overline{\mathcal{D}}_R=\left\{x \in \mathbb{R}^n, 0 \leq |x| \leq R/2\right\}$ and $y \in \mathcal{D}_R=\left\{y \in \mathbb{R}^n, 0 \leq |y| \leq R\right\}$, the triangle inequality gives $|x-y| \leq|x|+|y| \leq R/2+R \leq 2 R$. Consequently, $|x-y|^{n-\gamma} \leq 2^{n-\gamma} R^{n-\gamma}$ for all $x \in \overline{\mathcal{D}}_R,\,\, y \in \mathcal{D}_R$. This leads to the pointwise lower bound
$$
\mathcal{I}_{\gamma}\left(|u|^p\right)(t, x) \geq \frac{C_\gamma R^{-(n-\gamma)}}{2^{n-\gamma}} \int_{\mathcal{D}_R}|u(t, y)|^p d y \quad \text { for all } t \in (0, R^2), x \in \overline{\mathcal{D}}_R.
$$
Then, we get
\begin{align} \label{estimate_I_gamma}
&\int_0^{R^2} \int_{\overline{\mathcal{D}}_R} \mathcal{I}_{\gamma}\left(|u|^p\right)(t,x) \phi_R(x)\chi_R(t) d x d t \notag\\
\geq & \frac{R^{-(n-\gamma)}}{2^{n-\gamma}} \int_0^{R^2} \int_{\overline{\mathcal{D}}_R} \int_{\mathcal{D}_R} |u(t, y)|^p \phi_R(x)\chi_R(t) d y d x d t \notag\\
\geq & \operatorname{meas}(\overline{\mathcal{D}}_R) \frac{R^{-(n-\gamma)}}{2^{n-\gamma}} \int_0^{R^2} \int_{\mathcal{D}_R} |u(t, y)|^p \phi_R(y)\chi_R(t) d y d t \geq C R^\gamma \mathcal{J}_R,
\end{align}
where we have used the fact that $\phi_R \equiv 1$ on $\overline{\mathcal{D}}_R$, and $1 \geq \phi_R$ on $\mathcal{D}_R$. Moreover, a simple calculation leads to
\begin{equation} \label{estimate_initial_data}
\int_{0 \leq |x| \leq R}(u_0(x)+u_1(x))\phi_{R}(x)dx \geq C  \int_{0 \leq |x|\leq R}\langle x\rangle^{-n+q} \geq C  R^q.
\end{equation}
for $R\gg1$. Combining \eqref{I_u^p}, \eqref{estimate_I_gamma}, and \eqref{estimate_initial_data}, we infer that $$R^\gamma \mathcal{J}_R+ \varepsilon R^q \leq C \mathcal{J}_R^{\frac{1}{p}}R^{-2+\frac{n+2}{p'}}.$$ This leads to 
\begin{align*}
\varepsilon R^{q-\gamma} \leq C\mathcal{J}_R^{\frac{1}{p}}R^{-2-\gamma+\frac{n+2}{p'}} -\mathcal{J}_R \lesssim R^{-(2+\gamma)p'+n+2},
\end{align*}
that is 
\begin{align}
\varepsilon \lesssim R^{-(2+\gamma)p'+n+2+\gamma-q}. \label{eq1thr6.1}
\end{align}
The assumption \eqref{condition-Blowup} is equivalent to
\begin{align*}
-(2+\gamma)p'+n+2+\gamma-q < 0.
\end{align*}
We pass $R \to \infty$ in \eqref{eq1thr6.1} to derive a contradiction. This completes the proof of Theorem \ref{Thm_Blow_up}.
\end{proof}

\end{document}
